\documentclass[11pt]{article}
\usepackage{mathblatt}
\newcommand{\bereich}[1]{\par\addvspace{14pt}\noindent{\large\bfseries #1}\par\nobreak\vspace{4pt}}
\newcommand{\bl}[1]{\par\noindent #1\par\vspace{3pt}}
\newcommand{\blrand}[1]{\par\noindent{\color{mbgrau}#1}\par\vspace{3pt}}
\newcommand{\blhier}[1]{\par\noindent\llap{$\blacktriangleright$\hspace{0.5em}}\textbf{#1}\par\vspace{3pt}}
\newcommand{\stufe}[1]{\par\noindent\hspace*{1.5em}{\small #1}\par\vspace{2pt}}
\newcommand{\pkt}{\textperiodcentered{}}
\begin{document}

\pfheftstil
\pflaufkopf[Abitur GK \pkt{} Prüfung]{Baumdiagramm + bedingte Wahrscheinlichkeit}{AS20}\pfseitevon
\pfheftkopf{Baumdiagramm + bedingte Wahrscheinlichkeit}{}
\bl{Baumdiagramm}
\blhier{Ziehen ohne Zurücklegen}
\stufe{Ergebnisse aufzählen}
\stufe{Wahrscheinlichkeit gesucht}
\bl{Vierfeldertafel}
\bl{Bedingte Wahrscheinlichkeit}
\blrand{Term und Anteil am Baum}

\par\addvspace{14pt}\noindent Vorher: Geraden + Ebenen\par
\par\noindent Weiter: Binomialverteilung\par
\end{document}